% from ega4-17-fr.tex
\oldpage[IV-4]{73}

projections canoniques. Soit $V$ (resp. $V'$) l'ensemble des $x\in X$ (resp. $x'\in X'$) où $f$ possède l'une des propriétés suivantes (resp. où $f'$ possède la même propriété): être:
\begin{enumerate}
  \item[(i)] localement de type fini;
  \item[(ii)] localement de présentation finie;
  \item[(iii)] plat;
  \item[(iv)] non ramifié;
  \item[(v)] lisse;
  \item[(vi)] étale.
\end{enumerate}
Alors on a $V'=g'^{-1}(V)$ (autrement dit, pour que $f'$ ait la propriété considérée en un point $x'$, il faut et il suffit que $f$ ait cette propriété au point $x$).
\end{proposition}

La propriété (iii) n'est mise que pour mémoire, et ne nécessite pas l'hypothèse que $g$ soit localement de présentation finie ((2.2.11, (iv)), compte tenu du fait que la projection $g':X'\to X$ est un morphisme plat). En vertu de (17.7.2), les assertions relatives aux propriétés (iv), (v) et (vi) sont des conséquences de l'assertion relative à (ii). Il suffit donc de considérer les cas (i) et (ii). Il est clair que $V'\supset g'^{-1}(V)$; reste donc à montrer que $V'\subset g'^{-1}(V)$; notons que les ensembles $V$ et $V'$ sont ouverts, et $W=g'(V')$ est aussi ouvert dans $X$ en vertu de (2.4.6). Il s'agit donc de prouver que le morphisme $f|W:W\to Y$ est localement de type fini (resp. localement de présentation finie); comme par hypothèse le composé
\[
V'\xrightarrow{g''}W\xrightarrow{f|W}Y
\]
(où $g''$ est la restriction de $g'$), égal à $V'\xrightarrow{f'}Y'\xrightarrow{g}Y$, est localement de type fini (resp. localement de présentation finie) et que $g''$ est surjectif (donc fidèlement plat), on est ramené à prouver le lemme suivant, qui améliore (11.3.16):

\begin{lemma}[17.7.5]
\label{IV.17.7.5}
Soient $f:X\to Y$ un morphisme fidèlement plat et localement de présentation finie, $g:Y\to Z$ un morphisme tel que $g\circ f:X\to Z$ ait l'une des propriétés suivantes: être:
\begin{enumerate}
  \item[(i)] localement de type fini;
  \item[(ii)] localement de présentation finie;
  \item[(iii)] de type fini.
\end{enumerate}
Alors $g$ a la même propriété. Si en outre $f$ est quasi-compact ou $g$ quasi-séparé, la même conclusion est valable pour la propriété:
\begin{enumerate}
  \item[(iv)] être de présentation finie.
\end{enumerate}
\end{lemma}

Dans les cas (i) et (ii), il s'agit de voir que pour tout $y\in Y$, il y a un voisinage ouvert affine $V$ de $y$ dans $Y$ et un voisinage ouvert affine $W$ de $z=g(y)$ dans $Z$ contenant $g(V)$, tels que le morphisme $V\to W$, restriction de $g$, soit de type fini (resp. de présentation finie). Or, il existe par hypothèse un $x\in X$ tel que $f(x)=y$, un voisinage ouvert affine $U$ de $x$ dans $X$, un voisinage ouvert affine $V'$ de $y$ dans $Y$ contenant $f(U)$ et un voisinage ouvert affine $W$ de $z$ dans $Z$ contenant $g(V')$ tels que le morphisme $f_1:U\to V'$, restriction de $f$ soit plat et de présentation finie, et le morphisme $g_1:V'\to W$ restriction de $g$, tel que $g_1\circ f_1$ soit de type fini (resp. de présentation finie). Alors $f(U)$ est ouvert dans $Y$ (2.4.6), et si $V\subset f(U)$ est un voisinage affine de $y$, le morphisme $f_2:f_1^{-1}(V)\to V$,

